\section{Chapter summary}\label{section:disc-chapter-summary}

The five research questions and their results tell a connected story. At each stage, the pipeline gives up some silver agreement to achieve a property that the silver-label agreement does not reward: the grounding filter prioritizes source faithfulness, the design of the cascade adds provider diversity, and the provenance contract achieves auditability for each artifact and output. The common principle behind these decisions is that keeping an unsupported or untraceable claim in a clinical knowledge base is more harmful than dropping a claim that might be correct. The main negative result is also clear, and is stated plainly: the calibrated cascade is not cost-justified against a single strong model (single-Sonnet) under the defined silver metric (RQ3); however, that comparison is conservative, because an Opus silver-standard might structurally favor the Sonnet baseline, as they are both Anthropic models. 

The biggest contribution is therefore the framework itself, not the specific measurements. As discussed in Section~\ref{section:disc-agreement-based-cascading}, the reported cost-quality results depend on the corpus, the voter pool, and the model pricing at evaluation time; whereas the document-extraction sweep, the calibrated agreement cascade with its gates and the offline-replay protocol, and the NLI-based faithfulness checks are reusable and recalibratable. The threats discussed in Section~\ref{section:disc-threats-to-validity} bound these conclusions; Chapter~\ref{chapter:Conclusion_and_Future_Work} turns these findings into the final contributions of the thesis, and the work they leave open.